%!TEX root = ../main.tex
% =============================================================
%  CAPÍTULO 4 — MARCO TECNOLÓGICO       (Entrega 4 · 10 % · 25/8)
% =============================================================

\chapter{Marco tecnológico}\label{cap:marco-tecnologico}

\begin{guia}[Guía de la cátedra: qué se evalúa en este capítulo]
Documenta el \emph{cómo}. No es un manual de uso: es la justificación
de cada elección tecnológica. Regla: nada ``porque sí''.
\begin{enumerate}
    \item \textbf{Arquitectura}: diagrama con componentes y conexiones,
          responsabilidad y tecnología de cada uno, flujo de datos y
          justificación del estilo arquitectónico con alternativas
          evaluadas \citep{sommerville2016}.
    \item \textbf{Prototipo o \gls{g:mvp}}: cuál es y por qué; alcance
          implementado vs.\ simulado; capturas o enlace; cómo se probó.
    \item \textbf{Stack}: por tecnología, nombre, versión, capa, por qué
          (madurez, comunidad, costo, conocimiento del equipo) y qué
          alternativas se descartaron.
    \item \textbf{Infraestructura y base de datos}: dónde se aloja cada
          componente, escalabilidad, costos (enlazados al
          capítulo~\ref{cap:economica}), motor y modelo de datos.
    \item \textbf{Diseño y \gls{ux}}: quién lo usa, mockups, decisiones de
          diseño y \textbf{cómo se validó con usuarios} (la cátedra lo marca
          como importante; cuestionario en el apéndice~\ref{apx:encuesta}).
    \item \textbf{Seguridad, calidad y despliegue}: autenticación, control
          de acceso, cifrado, plan de pruebas, repositorio, \gls{cicd} y
          monitoreo.
\end{enumerate}
\end{guia}

\completar{Desarrollo del capítulo.}

\section{\completar{Nombre de la sección}}

\completar{Texto.}

\begin{figure}[H]
    \centering
    \includegraphics[width=0.85\textwidth]{example-image-a}
    \caption{Diagrama de arquitectura del sistema (ejemplo de figura)}
    \label{fig:arquitectura}
\end{figure}

\begin{ejemplo}[Ejemplo: decisión con alternativas]
Se adoptó una arquitectura \completar{cliente-servidor en capas} porque
\completar{el volumen de usuarios previsto y el equipo de una persona no
justifican la complejidad operativa de microservicios}. Se evaluó
\completar{alternativa} y se descartó porque \completar{motivo concreto}.
\end{ejemplo}

\section{\completar{Nombre de la sección}}

\completar{Texto.}

\begin{table}[H]
    \centering
    \caption{Stack tecnológico y alternativas evaluadas (ejemplo de estructura)}
    \label{tab:stack}
    \small
    \begin{tabularx}{\textwidth}{l l l L L}
        \toprule
        \tb{Capa} & \tb{Tecnología} & \tb{Versión} & \tb{Justificación} & \tb{Alternativas descartadas} \\
        \midrule
        Frontend & \completar{framework} & \completar{x.y} & \completar{madurez, comunidad, conocimiento del equipo} & \completar{alternativa: motivo} \\
        Backend & \completar{lenguaje/framework} & \completar{x.y} & \completar{...} & \completar{...} \\
        Base de datos & \completar{motor} & \completar{x.y} & \completar{relacional por integridad referencial y consultas ad hoc} \citep{postgresqldocs} & \completar{NoSQL: sin necesidad de esquema flexible} \\
        Modelo & \completar{librería} & \completar{x.y} & \completar{...} \citep{pedregosa2011} & \completar{...} \\
        \bottomrule
    \end{tabularx}
\end{table}

\begin{figure}[H]
    \centering
    \begin{subfigure}[b]{0.48\textwidth}
        \includegraphics[width=\textwidth]{example-image-a}
        \caption{Pantalla de inicio}
    \end{subfigure}
    \hfill
    \begin{subfigure}[b]{0.48\textwidth}
        \includegraphics[width=\textwidth]{example-image-b}
        \caption{Panel principal}
    \end{subfigure}
    \caption{Mockups de pantallas (ejemplo de figura compuesta)}
    \label{fig:mockups}
\end{figure}

\begin{lstlisting}[language=Python, caption={Ejemplo de código fuente con resaltado}, label={lst:ejemplo}]
def test_calcular_exposicion():
    # exposición = probabilidad x impacto
    assert calcular_exposicion(probabilidad=4, impacto=5) == 20
\end{lstlisting}
